\section{Main Hypothesis}\label{sec:attack:hypotheses}

Chapter~\ref{ch:hub} tested H1.
This chapter tests one further hypothesis, H2.
The second-round rewriting is reported as an exploratory result without a hypothesis number, and the remaining comparisons form the ablation study of Section~\ref{sec:attack:ablations}.

\paragraph{Hypothesis H2: matching improves reconstruction from a small leak.} The supervised fit trains the map on the rows of the leak alone.
The hybrid fit adds a matching term that reads the two unpaired pools and no text (Section~\ref{sec:attack:method}).
With a leak of \vf{90} rows, word-F1 is higher under the hybrid fit than under the supervised fit.
Criterion, written before the run: on Corpus A and on legal128w, the gain in word-F1 is at least \vf{0.10}, with the lower bound of its \vf{95} percent interval above \vf{0.05}. % src: w46/plan.md (prediction P1)

\paragraph{What H2 claims.} H2 is a claim about word-F1, the word overlap between a reconstruction and the true prompt, under the resources of Section~\ref{sec:attack:threat}.
It does not claim that prompts are recovered exactly, that their names and numbers are recovered, or that a stated amount of information leaks.
The results report exact match and decoded examples beside word-F1 so that the difference stays visible.

\paragraph{Exploratory question: second-round rewriting.} A language model receives the three decoders' reconstructions of a prompt and rewrites them into one text.
The question is whether word-F1 after the rewriting exceeds the mean word-F1 of the three reconstructions.
No criterion preceded the runs, the comparison with the best single decoder has no interval, and no control with mismatched prompts was run for this protocol, so Section~\ref{sec:attack:rewriting} reports the result as exploratory.
